% ---------------------------------------------------------------------------
% Experiment History -- subsection 9: Summary of Hypotheses and Verdicts.
% Sources: README A.0-A.7.
%
% Rewritten 2026-09-11: prose cut to four sentences (the tables are the
% content) and the verdict table cut from 17 rows to 13, with cell text
% shortened to plain wording.
%
% Two tables, two jobs:
%   tab:verdicts     -- what was asked, what answered it, what the answer
%                       was.  "Superseded by" is the column the supervisor
%                       reads: it shows retracted claims were retracted BY
%                       US, by a named later measurement.
%   tab:traceability -- convention (a): the prose names no simulation, so
%                       the mapping to directory and job ID lives here once.
%
% PLACEMENT: second-to-last.  "What the History Determines" stays the
% closing subsection -- a reference table is a bad last word for an appendix
% meant to hand off forward.
%
% The dashes in the job column are accurate: those runs predate job-ID
% logging in the run index.
% ---------------------------------------------------------------------------

\subsection{Summary of Hypotheses and Verdicts}
\label{subsec:exphist:verdicts}
Table~\ref{tab:verdicts} collects what each step was built to answer.
Four rows are verdicts against claims we had reported ourselves, and each names the later measurement
that overturned it.
The prose above names no simulation, since the code is not what the argument rests on, so
Table~\ref{tab:traceability} gives the mapping for anyone reproducing a number.

\begin{table*}[!t]
\renewcommand{\arraystretch}{1.3}
\caption{What each step asked, what answered it, and what the answer was.}
\label{tab:verdicts}
\centering
\begin{tabular}{p{0.29\textwidth} p{0.31\textwidth} l p{0.19\textwidth}}
\hline\hline
\textbf{Question} & \textbf{Evidence} & \textbf{Verdict} & \textbf{Superseded by} \\
\hline
A power threshold is evaded by power tuning alone
& BER $0.19$ at $10\,\%$ detection, power held just below the threshold
& yes & --- \\

A diagonal Gaussian policy can produce QPSK-like received statistics
& detection $100\,\%$, kurtosis stalled at $-0.25$
& no, structurally & --- \\

A more expressive action distribution repairs it under policy gradient
& flow and mixture heads stall at the same kurtosis, one run of $10^6$ steps
& no & --- \\

The same flow trained by direct gradient evades the kurtosis test
& kurtosis $-1.31$ at BER $0.170$, detection $\approx 0\,\%$
& yes, with a genie observation
& \S\ref{subsec:exphist:untrainable} \\

Two coordinated agents beat one at equal total power
& BER $0.195$ against $0.170$ at total power $\approx 0.97$
& yes, by $\approx 15\,\%$
& retracts our ``roughly double'' \\

Policy-gradient learning over raw IQ yields a stealthy jammer
& detector output pinned at $0.999$ in every run, entropy flat
& no, structurally
& perturbation parameters \\

Effectiveness and stealth trade off fundamentally, capping BER near $0.04$
& the cap was guard-bin occupancy, not a detector property
& no
& \S\ref{subsec:exphist:characterisation} \\

The CNN's $99.79\,\%$ measures resistance to jamming
& flags power $0.03$ out-of-band causing no errors, misses in-band at BER $0.36$
& no & --- \\

The in-band blind spot is a gap in the training data
& retraining closes it at accuracy $99.8 \to 90.5\,\%$, false alarms $0 \to 3.8\,\%$
& no, it is a trade-off & --- \\

A stealthy attacker reaching BER $0.42$ exists on the lossless channel
& against the suite, at most $0.005$ at any usable budget
& no
& energy detector \\

A sparse in-band jammer imposes a BER floor independent of SNR
& $0.05$--$0.07$ at eight subcarriers across 5--30\,dB
& yes & --- \\

Channel-aware subcarrier selection beats blind selection
& $+70\,\%$ at matched configuration, $-4.6$ to $+6.2\,\%$ at matched detectability
& no, once detectability is matched
& \S\ref{subsec:exphist:realistic} \\

An effective and unflagged attacker exists on a realistic channel
& at the defender's own false-alarm rate: none at 5--15\,dB, BER $0.011$ and $0.004$ at 20 and 30\,dB
& yes, and modest
& retracts our $\varepsilon \leq 0.5$ figure \\
\hline\hline
\end{tabular}
\end{table*}

\begin{table}[!t]
\renewcommand{\arraystretch}{1.3}
\caption{Narrative step to source directory and cluster job.}
\label{tab:traceability}
\centering
\begin{tabular}{l l l}
\hline\hline
\textbf{Step} & \textbf{Directory} & \textbf{Job(s)} \\
\hline
Chain validation          & \texttt{simulation00}, \texttt{04b} & --- \\
Power threshold           & \texttt{simulation01}  & --- \\
Kurtosis, Gaussian head   & \texttt{simulation02}  & --- \\
Kurtosis, flow, PPO       & \texttt{simulation03}  & 93396--93423 \\
Kurtosis, flow, direct    & \texttt{simulation03b} & 98432--98516 \\
Kurtosis, mixture head    & \texttt{simulation03c} & 98764--98777 \\
Two agents, direct        & \texttt{simulation04}  & 99211, 100041 \\
Detector, flat carrier    & \texttt{simulation05}  & 101306, 101314 \\
Detector, OFDM            & \texttt{simulation06}  & 101328 \\
Jammer, white-box         & \texttt{simulation06}  & --- \\
Jammer, 2-D action space  & \texttt{simulation06b} & --- \\
Jammer, black-box         & \texttt{simulation07}  & 101622--101817 \\
Characterisation          & \texttt{frontier}      & 101860, 101866 \\
Complex transform, suite  & \texttt{frontier}      & 102115 \\
Realistic channel         & \texttt{simulation08}  & 101870, 102305 \\
Channel-valid detector    & \texttt{simulation08}  & 102316, 102319 \\
Matched detectability     & \texttt{simulation08}  & 102390 \\
Minimal model             & \texttt{m0}            & 2243867, 2243879 \\
\hline\hline
\end{tabular}
\end{table}
